\begin{corollary}[$D_1$ and $P_1$ on the families]
\label{cor:c11-D1fam}
Under the hypotheses of Definition \ref{def:c10-hyp}, on the cohomology of the part with the fewest fields of every family of Definition \ref{def:c11-families}, with every free direction of Definition \ref{def:c11-free},
\[
D_1\cong4\,(\text{the map})-\sum_a\big(\theta_{X_a}+\tfrac13\big)\,\text{units}\cdot v.X_a,\qquad P_1\cong-\sum_a\big(\theta_{P,X_a}+\tfrac13\big)\,\text{units}\cdot\partial_bX_a,
\]
with the map of Theorem \ref{thm:c10-brst}, the position part of $D_1$ and the component $P_1$ taken along $b=n$, and along every $b\in V$ on twist two for $S'\le4$ and on the letters without derivatives (Theorem \ref{thm:c11-families}), where $\theta_{X_a}$ is $\theta$, $\theta_\lambda$ or $\theta_V$ and $\theta_{P,X_a}$ is $\theta_P$, $\theta_{P,\lambda}$ or $\theta_{P,V}$ according to the sector of the letter $X_a$, $v.X_a$ and $\partial_bX_a$ act on the letter $X_a$ alone, and the directions $Q\int(\iota_xA|_V,c^*)$, $Q\int(\rho(x)\lambda^*,\lambda^*)$ contribute $s$-exact terms. At the values $\theta=\theta_\lambda=\theta_V=-\tfrac13$ of Theorem \ref{thm:c9-theta2}, which twist two, the fermion and the gluon $\mathfrak{sl}(2)$ sectors reproduce under the hypotheses of Theorem \ref{thm:c11-arity2} (Proposition \ref{prop:c11-theta-sector}), $D_1$ induces four times the map of Theorem \ref{thm:c10-brst} on every family, with the spectra of Table \ref{tab:c11-results}, $-8$ units times the eigenvalues of $\sum_lH_{l,l+1}$, which is twice the one-loop anomalous dimensions; and $P_1\cong0$ in the translation-covariant scheme, while for equal coefficients $\theta_P$ along $P_b$, $P_1\cong(\theta-\theta_P)$ units times $\partial_b$ with $\theta=-\tfrac13$.
\end{corollary}
